\chapter{2012年陕西卷（理）}

\section{单选题}

\begin{problem}
集合 \(M = \{x \mid \lg x > 0\}\)，\(N = \{x \mid x^2 \leqslant 4\}\)，则 \(M \cap N =\)（\quad）

\choices
    {(1,2)}
    {\([1,2)\)}
    {(1,2]}
    {[1, 2]}
\end{problem}

\begin{problem}
下列函数中，既是奇函数又是增函数的为（\quad）

\choices
    {\(y = x + 1\)}
    {\( y = -x^3 \)}
    {\( y = \frac{1}{x} \)}
    {\(y = x|x|\)}
\end{problem}

\begin{problem}
设 \(a, b \in \R_+\) 是虚数单位，则“\(ab = 0\)”是“复数 \(a + \frac{b}{4}\) 为纯虚数”的（\quad）

\choices
    {充分不必要条件}
    {必要不充分条件}
    {充分必要条件}
    {既不充分也不必要条件}
\end{problem}

\begin{problem}
已知圆 \(C: x^{2} + y^{2} - 4x = 0\)，且过点 \(P(3,0)\) 的直线，则（\quad）

\choices
    {\(l\) 与 \(C\) 相交}
    {\(l\) 与 \(C\) 相切}
    {\(l\) 与 \(C\) 相离}
    {以上三个选项均有可能}
\end{problem}

\begin{problem}
如图，在空间直角坐标系中有直三棱柱 \(ABC - A_{1}B_{1}C_{1}\)，\(CA = CC_{1} = 2CB\)，则直线 \(BC_{1}\) 与直线 \(AB_{1}\) 夹角的余弦值为（\quad）

\choices
    {\(\frac{\sqrt{5}}{5}\)}
    {\( \frac{\sqrt{5}}{3} \)}
    {\( \frac{2\sqrt{5}}{5} \)}
    {\(\frac{3}{5}\)}
\end{problem}

\begin{problem}
从甲、乙两个城市分别随机抽取 16 台自动售货机. 对其销售额进行统计，统计数据用茎叶图表示 (如图所示). 设甲、乙两组数据的平均数分别为 \(\overline{x}\)，\(\overline{x}\)，中位数分别为 \(m_{\text{甲}}, m_{\text{乙}}\)，则（\quad）

\choices
    {\(\overline{x}_\text{甲} < \overline{x}_\text{乙},m_{\text{甲}} > m_\text{乙}\)}
    {\(\overline{x}_\text{甲} < \overline{x}_\text{乙},m_{\text{甲}} < m_\text{乙}\)}
    {\(\overline{x}_{\text{甲}} > \overline{x}_{\text{乙}}, m_{\text{甲}} > m_{\text{乙}}\)}
    {\(\overline{x}_{\text{甲}} > \overline{x}_{\text{乙}}, m_{\text{甲}} < m_{\text{乙}}\)}
\end{problem}

\begin{problem}
设函数 \( f(x) = x\e^{x} \)，则（\quad）

\choices
    {\(x = 1\) 为 \(f(x)\) 的极大值点}
    {\(x = 1\) 为 \(f(x)\) 的极小值点}
    {\(x = -1\) 为 \(f(x)\) 的极大值点}
    {\(x = -1\) 为 \(f(x)\) 的极小值点}
\end{problem}

\begin{problem}
两人进行乒乓球比赛，先赢 3 局者获胜，决出胜负为止，则所有可能出现的情形 (各入输赢局次的不同视为不同情形) 共有（\quad）

\choices
    {10 种}
    {15 种}
    {20 种}
    {30 种}
\end{problem}

\begin{problem}
在 \(\triangle ABC\) 中，角 \(A, B, C\) 所对的边长分别为 \(a, b, c\)，若 \(a^2 + b^2 = 2c^2\)，则 \(\cos C\) 的最小值为（\quad）

\choices
    {\( \frac{\sqrt{3}}{2} \)}
    {\( \frac{\sqrt{2}}{2} \)}
    {\( \frac{1}{2} \)}
    {\(-\frac{1}{2}\)}
\end{problem}

\begin{problem}
如图所示是用模板方法估计圆周率 \(\pi\) 值的程序框图，\(P\) 表示估计结果，则图中空白框内应填入（\quad）

\choices
    {\(P = \frac{N}{1000}\)}
    {\(P = \frac{4N}{1000}\)}
    {\(P = \frac{M}{1000}\)}
    {\(P = \frac{4M}{1000}\)}
\end{problem}

\section{填空题}

\begin{problem}
观察下列不等式: \[ 1 + \frac {1}{2 ^{2}} < \frac {3}{4}, \] \[ 1 + \frac {1}{2 ^{2}} + \frac {1}{4 ^{2}} < \frac {5}{3}, \] \[ 1 + \frac {1}{2 ^{2}} + \frac {2}{3 ^{2}} + \frac {3}{4 ^{2}} < \frac {7}{4}, \] 照此规律，第五个不等式为 \fillinblank{} \fillinblank{}.
\end{problem}

\begin{problem}
\((a + x)^{5}\) 展开式中 \(x^{2}\) 的系数为 10，则实数 \(a\) 的值为 \fillinblank{} \fillinblank{}.
\end{problem}

\begin{problem}
如图是抛物线形拱桥，当水面在 \(l\) 时，拱顶离水面 2 米，水面宽 4 米，水位下降 1 米后，水面宽 米 \fillinblank{}.
\end{problem}

\begin{problem}
设函数 \( f(x) = \begin{cases} \ln x, & x > 0, \\ -2x - 1, & x \leqslant 0, \end{cases} \) \( D \) 是由 \( z \) 轴和曲线 \( y = f(x) \) 及该曲线在点(1,0)处的切线所围成的封闭区域，则 \( x = -2x \) 与 \( D \) 上的最大值 为 \fillinblank{} \fillinblank{}.
\end{problem}

\begin{problem}
三选一. 【1】若存在实数 x 使 \( \left|x-a\right|+\left|x-1\right|\leqslant3 \) 成立，则实数 a 的取值范围是 \fillinblank{}. 【B】如图，在圆 O 中，直径 AB 与弦 CD 垂直，垂足为 E, EF \(\perp\) DB，垂足为 F，若 AB = 6, AE = 1，则 DF · DB =\fillinblank{}.

【C】直线 \(2y\cos \theta = 1\) 与圆 \(p = 2\cos \theta\) 相交的弦长为
\end{problem}

\section{解答题}

\begin{problem}
函数 \( f(x) = A \sin \left( \omega x - \frac{\pi}{6} \right) + 1 (A > 0, \omega > 0) \) 的最大值为 3，其图象相邻两条对称轴之间的距离为 \( \frac{\pi}{2} \). (1) 求函数 \( f(x) \) 的解析式. (2) 设 \(\alpha \in \left(0, \frac{\pi}{2}\right), f\left(\frac{\alpha}{2}\right) = 2\)，求 \(\alpha\) 的值.
\end{problem}

\begin{problem}
设 \(\{a_{n}\}\) 是公比不为1的等比数列，其前 \(n\) 项和为 \(S_{n}\)，且 \(a_{0}, a_{3}, a_{4}\) 成等差数列. (1) 求数列 \( \{a_{n}\} \) 的公比； (2) 证明: 对任意 \( k \in N_{+}, S_{k+2}, S_{k}, S_{k+1} \) 成等差数列.
\end{problem}

\begin{problem}
已知椭圆 \(C_1: \frac{x^2}{4} + y^2 = 1\)，椭圆 \(C_2\) 以 \(C_1\) 的长轴为短轴，且与 \(C_1\) 有相同的离心率. (1) 求椭圆 \(C_2\) 的方程; (2) 设 O 为坐标原点，点 A, B 分别在椭圆 \( C_{1} \) 和 \( C_{2} \) 上，\( \overrightarrow{OB} = 2\overrightarrow{OA} \) ，求直线 AB 的方程.
\end{problem}

\begin{problem}
设函数 \( f_{n}(x) = x^{n} + bx + c (n \in \mathbf{N}_{+}, b, c \in \R) \). (1) 设 \(n \geqslant 2, b = 1, c = -1\)，证明: \(f_n(x)\) 在区间 \(\left(\frac{1}{2}, 1\right)\) 内存在唯一零点; (2) 设 \(n = 2\)，若对任意 \(x_1, x_2 \in [-1, 1]\)，有 \(|f_2(x_1) - f_2(x_2)| \leqslant 4\)，求 \(b\) 的取值范围; (3) 在 (1) 的条件下，设 \(x_{n}\) 是 \(f_{n}(x)\) 在 \(\left(\frac{1}{2}, 1\right)\) 内的零点，判断数列 \(x_{2}, x_{3}, \cdots, x_{n}, \cdots\) 的增减性.
\end{problem}

\begin{problem}
(1) 如图，证明命题 “\(n\) 是平面 \(\pi\) 内的一条直线，\(b\) 是 \(\pi\) 外的一条直线 (\(b\) 不垂直于 \(\pi\)), \(c\) 是直线 \(b\) 在 \(\pi\) 上的投影，若 \(a \perp b\)，则 \(a \perp c\) “为真; (2) 写出上述命题的逆命题，并判断其真假（不需证明）.
\end{problem}

\begin{problem}
某银行柜台设有一个服务窗口，假设顾客办理业务所需的时间互相独立，且都是整数分钟，对以往顾客办理业务所需的时间统计结果如下: 办理业务所需的时间(分) 1 2 3 4 5 频率 0.1 0.4 0.3 0.1 0.1 从第一个顾客开始办理业务时计时. (1) 估计第三个顾客恰好等待 4 分钟开始办理业务的概率; (2) X 表示至第 2 分钟未已办理完业务的顾客人数，求 X 的分布列及数学期望.
\end{problem}

